\chapter{Implementacja systemu}
\label{cha:pierwszyDokument}

\section{Struktura projektu}

\begin{lstlisting}[style=styleBash, caption={Struktura katalogów projektu mnist-translator}, label=lst:struktura, numbers=none]
mnist-translator/
|-- main.py                  # Punkt wejscia CLI (brak logiki biznesowej)
|-- pyproject.toml           # Zarzadzanie zaleznosci (uv)
|-- src/
|   |-- data/
|   |   |-- ingest.py        # Ladowanie CSV i normalizacja pikseli
|   |   |-- extract.py       # Budowa macierzy cech X i etykiet y
|   |   |-- split.py         # Stratyfikowany podzial 80/10/10
|   |   `-- protocol.py      # Protokol DatasetSource
|   |-- model/
|   |   |-- train.py         # Trenowanie + zapis artefaktu
|   |   |-- evaluate.py      # Obliczanie metryk dokladnosci
|   |   `-- classifiers/
|   |       |-- svc.py       # Klasyfikator SVC z jadrem RBF
|   |       |-- random_forest.py  # Las Losowy (200 drzew)
|   |       `-- cnn.py       # Konwolucyjna siec neuronowa
|   `-- inference/
|       |-- camera.py        # Przechwytywanie wideo + MediaPipe
|       |-- csv_source.py    # Zrodlo cech z wiersza CSV
|       |-- predict.py       # Ladowanie modelu i predykcja
|       `-- protocol.py      # Protokol FeatureSource
|-- tests/                   # Testy jednostkowe (pytest)
|-- diagrams/                # Diagram architektury
`-- artifacts/               # Wytrenowane modele (gitignored)
\end{lstlisting}

\section{Interfejs wiersza poleceń}

Plik \texttt{main.py} udostępnia trzy tryby działania:

\begin{table}[H]
    \centering
    \caption{Tryby działania aplikacji mnist-translator}
    \label{tab:tryby}
    \begin{tabularx}{\textwidth}{llX}
        \toprule
        \textbf{Tryb} & \textbf{Opis} & \textbf{Kluczowe opcje} \\
        \midrule
        \texttt{train}  & Trenuje klasyfikator i zapisuje artefakt
                        & \texttt{-{}-classifier \{svc,rf,cnn\}}, \texttt{-{}-model}, \texttt{-{}-dataset} \\
        \texttt{eval}   & Ewaluuje model na zbiorze testowym
                        & \texttt{-{}-model}, \texttt{-{}-dataset} \\
        \texttt{infer}  & Inferencja z CSV lub z kamery
                        & \texttt{-{}-source \{csv,camera\}}, \texttt{-{}-row}, \texttt{-{}-camera} \\
        \bottomrule
    \end{tabularx}
\end{table}

\newpage
Przykładowe użycie:
\begin{lstlisting}[style=stylePython,caption={Przyklady uruchomien CLI}]
uv run python main.py train --classifier svc --model artifacts/svc_model.pkl
uv run python main.py eval --model artifacts/svc_model.pkl --dataset data/raw/sign_mnist_test.csv
uv run python main.py infer --source csv --row 42 --dataset data/raw/sign_mnist_test.csv
uv run python main.py infer --source camera --camera 0 --model artifacts/svc_model.pkl
\end{lstlisting}

\section{Implementacja klasyfikatorów}

\subsection{SVC z jądrem RBF}

\begin{lstlisting}[style=stylePython, caption={Budowa potoku SVC (src/model/classifiers/svc.py)}, label=lst:svc]
from sklearn.svm import SVC
from sklearn.pipeline import Pipeline
from sklearn.preprocessing import StandardScaler

def build_svc() -> Pipeline:
    """Buduje potok: standaryzacja + SVC z jadrem RBF."""
    return Pipeline([
        ("scaler", StandardScaler()),
        ("svc", SVC(kernel="rbf", C=10, gamma="scale",
                    decision_function_shape="ovr")),
    ])
\end{lstlisting}

\subsection{Las Losowy}

\begin{lstlisting}[style=stylePython, caption={Budowa Lasu Losowego (src/model/classifiers/random\_forest.py)}, label=lst:rf]
from sklearn.ensemble import RandomForestClassifier

def build_random_forest() -> RandomForestClassifier:
    """200 drzew, tryb OOB wlaczony."""
    return RandomForestClassifier(
        n_estimators=200,
        random_state=42,
        oob_score=True,
        n_jobs=-1,
    )
\end{lstlisting}

\newpage
\subsection{Konwolucyjna sieć neuronowa}

\begin{lstlisting}[style=stylePython, caption={Architektura CNN (src/model/classifiers/cnn.py)}, label=lst:cnn]
import tensorflow as tf

def build_cnn(num_classes: int = 24) -> tf.keras.Model:
    """Dwa bloki Conv2D->MaxPool + warstwa wyjsciowa z softmax."""
    model = tf.keras.Sequential([
        tf.keras.layers.Input(shape=(28, 28, 1)),
        tf.keras.layers.Conv2D(32, (3, 3), activation="relu",
                               padding="same"),
        tf.keras.layers.MaxPooling2D(),
        tf.keras.layers.Conv2D(64, (3, 3), activation="relu",
                               padding="same"),
        tf.keras.layers.MaxPooling2D(),
        tf.keras.layers.Flatten(),
        tf.keras.layers.Dense(128, activation="relu"),
        tf.keras.layers.Dropout(0.3),
        tf.keras.layers.Dense(num_classes, activation="softmax"),
    ])
    model.compile(optimizer="adam",
                  loss="sparse_categorical_crossentropy",
                  metrics=["accuracy"])
    return model
\end{lstlisting}

\section{Serializacja modeli}

Modele scikit-learn są zapisywane i wczytywane za pomocą biblioteki \texttt{joblib}:

\begin{lstlisting}[style=stylePython, caption={Zapis i odczyt modelu scikit-learn}, label=lst:joblib]
import joblib

# Zapis modelu
joblib.dump(model, "artifacts/svc_model.pkl")

# Wczytanie modelu
model = joblib.load("artifacts/svc_model.pkl")
\end{lstlisting}

Model CNN jest zapisywany w formacie natywnym Keras (rozszerzenie \texttt{.keras}):

\begin{lstlisting}[style=stylePython, caption={Zapis i odczyt modelu CNN (Keras)}, label=lst:keras]
# Zapis
model.save("artifacts/cnn_model.keras")

# Wczytanie
import tensorflow as tf
model = tf.keras.models.load_model("artifacts/cnn_model.keras")
\end{lstlisting}

\section{Główne zależności}

\begin{table}[H]
    \centering
    \caption{Główne zależności projektu (pyproject.toml)}
    \label{tab:deps}
    \begin{tabular}{ll}
        \toprule
        \textbf{Biblioteka} & \textbf{Zastosowanie} \\
        \midrule
        \texttt{numpy}, \texttt{pandas}   & Manipulacja danymi i macierzami \\
        \texttt{scikit-learn}             & SVC, Random Forest, metryki \\
        \texttt{joblib}                   & Serializacja modeli \\
        \texttt{opencv-python}            & Przechwytywanie klatek z kamery \\
        \texttt{mediapipe}                & Detekcja dłoni i punktów kluczowych \\
        \texttt{tensorflow >= 2.21.0}     & Budowa i trenowanie CNN \\
        \bottomrule
    \end{tabular}
\end{table}

\section{Architektura logiczna}
\label{sec:architektura_logiczna}

Projekt został podzielony na trzy warstwy:
\begin{itemize}
    \item \texttt{src/data} -- przygotowanie danych i podziały zbioru,
    \item \texttt{src/model} -- uczenie i ewaluacja klasyfikatorów,
    \item \texttt{src/inference} -- predykcja offline i online (kamera).
\end{itemize}

Warstwa wejściowa CLI znajduje się w \texttt{main.py} i deleguje logikę do
modułów w \texttt{src/}. Takie podejście upraszcza testowanie i ogranicza
sprzężenie między interfejsem użytkownika a logiką przetwarzania.

\section{Przetwarzanie danych}
\label{sec:przetwarzanie_danych}

Klasa \texttt{CsvDataset} ładuje plik CSV i wymusza spójność schematu
(obecność kolumny \texttt{label} oraz dokładnie 784 cechy obrazu). Dla każdej
próbki zwracana jest para:
\begin{center}
    \texttt{(feature\_vector: np.ndarray, label: str)}
\end{center}

Wektor cech jest typu \texttt{float32} i przyjmuje wartości w $[0,1]$.
Etykiety numeryczne przekształcane są do liter poprzez mapowanie:
\begin{equation}
    y_{letter} = \mathrm{chr}(\mathrm{ord}('A') + y_{int})
\end{equation}

Budowa macierzy cech realizowana jest przez \texttt{build\_feature\_matrix},
a podział danych przez \texttt{stratified\_split}. Domyślnie stosowany jest
podział 80/10/10 z zachowaniem proporcji klas.

\section{Uczenie modeli}
\label{sec:uczenie_modeli}

Moduł \texttt{src/model/train.py} obsługuje trzy warianty klasyfikatorów:
\begin{itemize}
    \item \textbf{SVC (RBF)} -- domyślny klasyfikator oparty o jądro radialne,
    \item \textbf{Random Forest} -- las losowy o 200 drzewach,
    \item \textbf{CNN (Keras)} -- sieć konwolucyjna dla wejścia 28$\times$28.
\end{itemize}

Modele klasyczne serializowane są przez \texttt{joblib} do formatu
\texttt{.pkl}. Model konwolucyjny zapisywany jest jako \texttt{.keras}
oraz osobny plik mapowania etykiet \texttt{.labels.json}.

\section{Ewaluacja jakości}
\label{sec:ewaluacja_jakosci}

Funkcja \texttt{evaluate} oblicza dokładność globalną oraz raport klasyfikacji
\texttt{classification\_report} biblioteki scikit-learn \cite{sklearn_docs}.
Wyniki zwracane są również jako słownik, co upraszcza dalsze automatyczne
raportowanie.

\section{Inferencja: CSV i kamera}
\label{sec:inferencja_csv_kamera}

Inferencja offline z jednego wiersza CSV realizowana jest przez
\texttt{CsvRowSource}. Inferencja online wykorzystuje klasę \texttt{CameraSource},
która:
\begin{enumerate}
    \item pobiera klatkę z kamery (OpenCV \cite{opencv_docs}),
    \item wykrywa dłoń przez MediaPipe GestureRecognizer \cite{mediapipe_docs},
    \item wyznacza obszar dłoni i dokonuje kadrowania,
    \item przeskalowuje obraz do 28$\times$28 oraz normalizuje piksele,
    \item przekazuje wektor 784 cech do modelu.
\end{enumerate}

W przypadku braku detekcji dłoni moduł zwraca wektor zerowy, dzięki czemu
interfejs inferencji pozostaje stabilny.

\section{Rola MediaPipe w systemie}

Kluczową rolą MediaPipe jest dostarczenie bounding boxa dłoni, dzięki czemu
przycięty i przeskalowany fragment klatki jest znacznie bardziej zbliżony
do próbek treningowych niż surowy obraz z kamery. Bez tego kroku predykcje
byłyby praktycznie bezużyteczne. Jednocześnie MediaPipe bywa zawodny
przy słabym oświetleniu lub gdy dłoń nie jest wyraźnie wyeksponowana.